\subsection{مبانی نظری و فرمول‌بندی ریاضی}
مسئله بهینه‌سازی هرس مجاورتی برای ماتریس وزن اولیه $\mathbf{W}_0$ به صورت ترکیب کمینه‌سازی خطای بازسازی فروبنیوس با جریمه تنک‌سازی و قید هندسی ساختاریافته تعریف می‌شود:
\begin{equation}
\min_{\mathbf{W}} \frac{1}{2} \|\mathbf{W} - \mathbf{W}_0\|_F^2 + \lambda \|\mathbf{W}\|_1 + \mathcal{I}_{\mathcal{S}}(\mathbf{W})
\end{equation}
که در آن $\lambda > 0$ ضریب منظم‌سازی مجاورتی، $\|\mathbf{W}\|_1 = \sum_{i,j} |W_{ij}|$ نرم $\ell_1$ جهت ترویج تنکی و $\mathcal{I}_{\mathcal{S}}(\mathbf{W})$ تابع نشانگر مجموعه مقید تنک $\mathcal{S}$ است:
\begin{equation}
\mathcal{I}_{\mathcal{S}}(\mathbf{W}) = \begin{cases} 0 & \mathbf{W} \in \mathcal{S} \\ +\infty & \mathbf{W} \notin \mathcal{S} \end{cases}
\end{equation}

الگوریتم حل از طریق عملگرهای تکراری زیر پیاده‌سازی می‌گردد:
\begin{enumerate}
    \item \textbf{عملگر آستانه‌گذاری نرم (\lr{Soft-Thresholding Operator}):}
    عملگر مجاورتی نرم $\ell_1$ جهت صفرسازی مقادیر ناچیز و انقباض وزن‌ها اعمال می‌شود:
    \begin{equation}
    \mathcal{S}_{\lambda}(\mathbf{W}) = \operatorname{sign}(\mathbf{W}) \odot \max(|\mathbf{W}| - \lambda, 0)
    \end{equation}

    \item \textbf{پروجکشن روی مجموعه مقید تنکی (\lr{Sparsity Projection}):}
    تصویر ماتریس روی ساختار نیمه‌ساختاریافته انویدیا $\Pi_{2:4}$ (حفظ ۲ عنصر با بیشترین قدرمطلق در هر بلوک ۴ تایی) یا تصویر بلوکی/صدکی $\Pi_{\mathcal{S}}$ انجام می‌پذیرد:
    \begin{equation}
    \Pi_{\mathcal{S}}(\mathbf{W}) = \operatorname{Top-}(1-s)(\mathbf{W}) \quad \text{یا} \quad \Pi_{2:4}(\mathbf{W})
    \end{equation}

    \item \textbf{تکرارهای گرادیان مجاورتی متناوب:}
    در هر گام تکرار $t \in \{1, \dots, T\}$ با نرخ گام $\eta$:
    \begin{equation}
    \mathbf{W}^{(t+1)} = \mathbf{W}^{(t)} - \eta \left( \mathbf{W}^{(t)} - \Pi_{\mathcal{S}}(\mathcal{S}_{\lambda}(\mathbf{W}^{(t)})) \right)
    \end{equation}
    در پایان، ماسک باینری دائمی $\mathbf{M} \in \{0, 1\}^{d_{\text{out}} \times d_{\text{in}}}$ تشکیل شده و وزن منجمد نهایی با رابطه $\mathbf{W}_{\text{final}} = \mathbf{W}_0 \odot \mathbf{M}$ تولید می‌شود.
\end{enumerate}


\paragraph{تحلیل طیفی و نامساوی‌های پایداری خطا}
برای بررسی میزان دقت تقریب ماتریس وزن یا نگاشت توجه در ماژول \lr{proxsparse}، تفاضل ماتریس فاکتورگیری‌شده از نگاشت متراکم اصلی بر اساس نامساوی هافدینگ-فون‌نیومان ارزیابی می‌گردد. فرض کنید $\mathbf{W} \in \mathbb{R}^{d_{\text{in}} \times d_{\text{out}}}$ ماتریس وزن هدف باشد که توسط تقریب کم‌مرتبه یا پراکنده با رتبه موثر $r$ جایگزین شده است. کران بالای خطای بازسازی بر مبنای نرم فروبنیوس به صورت زیر تصریح می‌شود:
\begin{equation}
\label{eq:spectral_bound_proxsparse}
\| \mathbf{W} - \mathbf{W}_{r} \|_F = \sqrt{\sum_{i=r+1}^{\min(d_{\text{in}}, d_{\text{out}})} \sigma_i^2(\mathbf{W})} \le \epsilon \| \mathbf{W} \|_F
\end{equation}
که در آن $\sigma_i(\mathbf{W})$ بیانگر $i$-امین مقدار منفرد مرتب‌شده ماتریس به صورت نزولی است. نسبت حفظ انرژی طیفی $\eta_r = \frac{\sum_{i=1}^r \sigma_i^2}{\sum_{i=1}^{d} \sigma_i^2}$ تضمین می‌کند که با انتخاب آستانه $\tau \ge 0.90$، بیش از نود درصد از پراکندگی واریانس داده‌ها حفظ شده و افت سیگنال در انتشار گرادیان رو به عقب به حداقل ممکن کاهش می‌یابد.

\paragraph{تحلیل مجانبی و پیچیدگی اشتراک حافظه سخت‌افزاری}
از منظر پیچیدگی محاسباتی و انتقال داده‌ها در سلسومراتب حافظه (حافظه سراسری \lr{HBM} به حافظه محلی \lr{SRAM})، عملیات ضرب ماتریسی در این مکانیزم نیازمند تخصیص بافر موقت با ابعاد فاکتورگیری‌شده است. برای پردازش دنباله‌ای به طول $N$ با بعد مخفی $D$، بار محاسباتی هسته از فرم درجه دو به رابطه زیر تقلیل می‌یابد:
\begin{equation}
\label{eq:complexity_proxsparse}
\mathcal{C}_{\text{proxsparse}} = \mathcal{O}\left( B \cdot N \cdot D \cdot r + B \cdot H \cdot N \cdot d_k \right) \ll \mathcal{O}\left( B \cdot H \cdot N^2 \cdot d_k \right)
\end{equation}
این بهبود مرتبه پیچیدگی، با کاهش چشمگیر حجم تبادل تانسورها در گذرگاه‌های ارتباطی کارت گرافیک، مانع از اشباع پهنای باند حافظه (پدیده \lr{Memory-Bound}) در فرآیند استنتاج طولانی‌مدت می‌گردد.
